% GENERATED FILE. Edit canonical lessons, then run npm run generate:books.
% canonical-source-hash: fnv1a64:f0f8534248cb052e
% canonical-lessons: KA-C509-nara, KA-C509-pakkelubu, KA-C509-bennumule, KA-C509-tolu, KA-C509-kuntu

\chapter{The Frame of the Body}
\label{ch:ka-the-frame-of-the-body}
\hlchaptermodality{509}

\begin{chapteropening}
\textbf{By the end of this chapter:} \emph{I can say a nerve, a rib, the spine, an arm and to limp.}

Complete the last lesson of chapter 509: I can say a nerve, a rib, the spine, an arm and to limp.
\end{chapteropening}

\section[\texorpdfstring{\kn{ನರ} (nara)}{ನರ (nara)}]{\kn{ನರ} (nara) — a nerve; a sinew}

\label{lesson:KA-C509-nara}



\begin{quote}
Before the new one: say the Kannada for the intestine, then the Kannada for a muscle.
\end{quote}

\subsection*{You'll want to know: \kn{ನರ}}
\textbf{\kn{ನರ}} — \emph{nara} — \textquotedblleft{}a nerve; a sinew\textquotedblright{}.

\begin{quote}
\kn{ವೈದ್ಯರು} \kn{ಅವನ} \kn{ಕೈಯ} \kn{ನರವನ್ನು} \kn{ಪರೀಕ್ಷಿಸಿದರು}.
\end{quote}

\begin{quote}
\emph{vaidyaru avana kaiya naravannu parīkṣisidaru.}
\end{quote}

\begin{quote}
The doctor examined the nerve in his arm.
\end{quote}

Plural: \kn{ನರಗಳು}, \emph{naragaḷu}.

In everyday talk also a vein; \kn{ನರಗಳು}, the nerves.

\begin{cousinweb}
\textbf{\kn{ನರ}} (\emph{nara}) is native Dravidian: DEDR 2903 lists Kannada \emph{nara}, \emph{naravu}, \textquotedblleft{}sinew, tendon, nerve, vein\textquotedblright{}.
\end{cousinweb}

\begin{grammarlens}[title={What you've built}]
One more word for describing what you lived and felt.
\end{grammarlens}

\subsection*{Your turn}
\begin{itemize}
\raggedright
  \item \emph{Say it:} \emph{nara}
  \item \emph{Say it:} \emph{nara}, once more
  \item \emph{Say it:} \emph{vaidyaru avana kaiya naravannu parīkṣisidaru.} — The doctor examined the nerve in his arm.
  \item \emph{Recall:} say the Kannada for the intestine, then the Kannada for a muscle, then say \emph{nara} again
\end{itemize}

\begin{tcolorbox}[breakable,colback=teal!4,colframe=teal!35!black,title={Before you move on}]
What does \kn{ನರ} mean? (\textquotedblleft{}A nerve; a sinew\textquotedblright{}.) Say it once more.
\end{tcolorbox}
\section[\texorpdfstring{\kn{ಪಕ್ಕೆಲುಬು} (pakkelubu)}{ಪಕ್ಕೆಲುಬು (pakkelubu)}]{\kn{ಪಕ್ಕೆಲುಬು} (pakkelubu) — a rib}

\label{lesson:KA-C509-pakkelubu}



\begin{quote}
Before the new one: say the Kannada for a muscle, then the Kannada for a nerve.
\end{quote}

\subsection*{You'll want to know: \kn{ಪಕ್ಕೆಲುಬು}}
\textbf{\kn{ಪಕ್ಕೆಲುಬು}} — \emph{pakkelubu} — \textquotedblleft{}a rib\textquotedblright{}.

\begin{quote}
\kn{ಬಿದ್ದಾಗ} \kn{ಅವನ} \kn{ಪಕ್ಕೆಲುಬು} \kn{ಮುರಿಯಿತು}.
\end{quote}

\begin{quote}
\emph{biddāga avana pakkelubu muriyitu.}
\end{quote}

\begin{quote}
When he fell, he broke a rib.
\end{quote}

\kn{ಪಕ್ಕ}, the side + \kn{ಎಲುಬು}, a bone: a bone of the side.

\begin{cousinweb}
\textbf{\kn{ಪಕ್ಕೆಲುಬು}} (\emph{pakkelubu}) is a Kannada compound: \emph{pakka}, \textquotedblleft{}the side\textquotedblright{}, which goes back to Sanskrit \emph{pakṣa}, \textquotedblleft{}wing, side\textquotedblright{}, through Prakrit \emph{pakkha} (CDIAL 7627), and \emph{elubu}, \textquotedblleft{}a bone\textquotedblright{}, which DEDR 839 lists among the native words for bone.
\end{cousinweb}

\begin{grammarlens}[title={What you've built}]
One more word for describing what you lived and felt.
\end{grammarlens}

\subsection*{Your turn}
\begin{itemize}
\raggedright
  \item \emph{Say it:} \emph{pakkelubu}
  \item \emph{Say it:} \emph{pakkelubu}, once more
  \item \emph{Say it:} \emph{biddāga avana pakkelubu muriyitu.} — When he fell, he broke a rib.
  \item \emph{Recall:} say the Kannada for a muscle, then the Kannada for a nerve, then say \emph{pakkelubu} again
\end{itemize}

\begin{tcolorbox}[breakable,colback=teal!4,colframe=teal!35!black,title={Before you move on}]
What does \kn{ಪಕ್ಕೆಲುಬು} mean? (\textquotedblleft{}A rib\textquotedblright{}.) Say it once more.
\end{tcolorbox}
\section[\texorpdfstring{\kn{ಬೆನ್ನುಮೂಳೆ} (bennumūḷe)}{ಬೆನ್ನುಮೂಳೆ (bennumūḷe)}]{\kn{ಬೆನ್ನುಮೂಳೆ} (bennumūḷe) — the spine; the backbone}

\label{lesson:KA-C509-bennumule}



\begin{quote}
Before the new one: say the Kannada for a nerve, then the Kannada for a rib.
\end{quote}

\subsection*{You'll want to know: \kn{ಬೆನ್ನುಮೂಳೆ}}
\textbf{\kn{ಬೆನ್ನುಮೂಳೆ}} — \emph{bennumūḷe} — \textquotedblleft{}the spine; the backbone\textquotedblright{}.

\begin{quote}
\kn{ಸರಿಯಾಗಿ} \kn{ಕುಳಿತರೆ} \kn{ಬೆನ್ನುಮೂಳೆಗೆ} \kn{ಒಳ್ಳೆಯದು}.
\end{quote}

\begin{quote}
\emph{sariyāgi kuḷitare bennumūḷege oḷḷeyadu.}
\end{quote}

\begin{quote}
Sitting properly is good for the spine.
\end{quote}

\kn{ಬೆನ್ನು}, the back + \kn{ಮೂಳೆ}, a bone.

\begin{cousinweb}
\textbf{\kn{ಬೆನ್ನುಮೂಳೆ}} (\emph{bennumūḷe}) is a compound of two native Dravidian words: \emph{bennu}, \textquotedblleft{}the back\textquotedblright{} (DEDR 5488), and \emph{mūḷe}, \textquotedblleft{}a bone\textquotedblright{} (DEDR 5050).
\end{cousinweb}

\begin{grammarlens}[title={What you've built}]
One more word for describing what you lived and felt.
\end{grammarlens}

\subsection*{Your turn}
\begin{itemize}
\raggedright
  \item \emph{Say it:} \emph{bennumūḷe}
  \item \emph{Say it:} \emph{bennumūḷe}, once more
  \item \emph{Say it:} \emph{sariyāgi kuḷitare bennumūḷege oḷḷeyadu.} — Sitting properly is good for the spine.
  \item \emph{Recall:} say the Kannada for a nerve, then the Kannada for a rib, then say \emph{bennumūḷe} again
\end{itemize}

\begin{tcolorbox}[breakable,colback=teal!4,colframe=teal!35!black,title={Before you move on}]
What does \kn{ಬೆನ್ನುಮೂಳೆ} mean? (\textquotedblleft{}The spine; the backbone\textquotedblright{}.) Say it once more.
\end{tcolorbox}
\section[\texorpdfstring{\kn{ತೋಳು} (tōḷu)}{ತೋಳು (tōḷu)}]{\kn{ತೋಳು} (tōḷu) — an arm; the upper arm}

\label{lesson:KA-C509-tolu}



\begin{quote}
Before the new one: say the Kannada for a rib, then the Kannada for the spine.
\end{quote}

\subsection*{You'll want to know: \kn{ತೋಳು}}
\textbf{\kn{ತೋಳು}} — \emph{tōḷu} — \textquotedblleft{}an arm; the upper arm\textquotedblright{}.

\begin{quote}
\kn{ಅಪ್ಪ} \kn{ಮಗುವನ್ನು} \kn{ತೋಳಿನಲ್ಲಿ} \kn{ಎತ್ತಿಕೊಂಡರು}.
\end{quote}

\begin{quote}
\emph{appa maguvannu tōḷinalli ettikoṇḍaru.}
\end{quote}

\begin{quote}
Father picked the baby up in his arms.
\end{quote}

Plural: \kn{ತೋಳುಗಳು}, \emph{tōḷugaḷu}.

\kn{ಕೈ} is the hand or the whole arm; \kn{ತೋಳು} is from shoulder to elbow, and also a sleeve.

\begin{cousinweb}
\textbf{\kn{ತೋಳು}} (\emph{tōḷu}) is native Dravidian: DEDR 3564 lists Kannada \emph{tōḷ(u)}, \textquotedblleft{}the arm\textquotedblright{}, beside Tamil \emph{tōḷ}, \textquotedblleft{}shoulder, arm\textquotedblright{}.
\end{cousinweb}

\begin{grammarlens}[title={What you've built}]
One more word for describing what you lived and felt.
\end{grammarlens}

\subsection*{Your turn}
\begin{itemize}
\raggedright
  \item \emph{Say it:} \emph{tōḷu}
  \item \emph{Say it:} \emph{tōḷu}, once more
  \item \emph{Say it:} \emph{appa maguvannu tōḷinalli ettikoṇḍaru.} — Father picked the baby up in his arms.
  \item \emph{Recall:} say the Kannada for a rib, then the Kannada for the spine, then say \emph{tōḷu} again
\end{itemize}

\begin{tcolorbox}[breakable,colback=teal!4,colframe=teal!35!black,title={Before you move on}]
What does \kn{ತೋಳು} mean? (\textquotedblleft{}An arm; the upper arm\textquotedblright{}.) Say it once more.
\end{tcolorbox}
\section[\texorpdfstring{\kn{ಕುಂಟು} (kuṇṭu)}{ಕುಂಟು (kuṇṭu)}]{\kn{ಕುಂಟು} (kuṇṭu) — to limp}

\label{lesson:KA-C509-kuntu}



\begin{quote}
Before the new one: say the Kannada for the spine, then the Kannada for an arm.
\end{quote}

\subsection*{You'll want to know: \kn{ಕುಂಟು}}
\textbf{\kn{ಕುಂಟು}} — \emph{kuṇṭu} — \textquotedblleft{}to limp\textquotedblright{}.

\begin{quote}
\kn{ಕಾಲಿಗೆ} \kn{ಗಾಯ} \kn{ಆಗಿದ್ದರಿಂದ} \kn{ಅವನು} \kn{ಕುಂಟುತ್ತಿದ್ದನು}.
\end{quote}

\begin{quote}
\emph{kālige gāya āgiddarinda avanu kuṇṭuttiddanu.}
\end{quote}

\begin{quote}
Because his leg was hurt, he was limping.
\end{quote}

Past (I …): \kn{ಕುಂಟಿದೆ}, \emph{kuṇṭide}.

\kn{ಕುಂಟ} is a lame man; \kn{ಕುಂಟಾಟ} is hopscotch.

\begin{cousinweb}
\textbf{\kn{ಕುಂಟು}} (\emph{kuṇṭu}) is native Dravidian: DEDR 1688 lists Kannada \emph{kuṇṭu}, \textquotedblleft{}to limp, halt, hop\textquotedblright{}, and \emph{kuṇṭa}, \textquotedblleft{}a lame man\textquotedblright{}.
\end{cousinweb}

\begin{grammarlens}[title={What you've built}]
One more word for describing what you lived and felt.
\end{grammarlens}

\subsection*{Your turn}
\begin{itemize}
\raggedright
  \item \emph{Say it:} \emph{kuṇṭu}
  \item \emph{Say it:} \emph{kuṇṭu}, once more
  \item \emph{Say it:} \emph{kālige gāya āgiddarinda avanu kuṇṭuttiddanu.} — Because his leg was hurt, he was limping.
  \item \emph{Recall:} say the Kannada for the spine, then the Kannada for an arm, then say \emph{kuṇṭu} again
\end{itemize}

\begin{tcolorbox}[breakable,colback=teal!4,colframe=teal!35!black,title={Before you move on}]
What does \kn{ಕುಂಟು} mean? (\textquotedblleft{}To limp\textquotedblright{}.) Say it once more.
\end{tcolorbox}
